\documentclass[11pt]{article}
\usepackage[margin=1in]{geometry}
\usepackage{amsmath,amssymb,amsthm}
\usepackage[hidelinks]{hyperref}

\newtheorem{theorem}{Theorem}
\newtheorem{corollary}[theorem]{Corollary}
\newtheorem{lemma}[theorem]{Lemma}
\newcommand{\R}{\mathbb R}
\newcommand{\Z}{\mathbb Z}

\title{Three cyclic chord classes force odd regularity\\
\large A joint-realizability obstruction for Erd\H{o}s Problem 132}
\author{Blake Stone}
\date{27 November 2026}

\begin{document}
\maketitle

\begin{abstract}
We prove an elementary rigidity lemma for an odd cyclic labeling of distinct
planar points.  If the label differences one, two, and three each determine
a single Euclidean length, then the points form a regular star polygon in
label order.  Consequently, a natural family of complete-graph colorings
with exactly one ``rare'' color cannot be realized as the equality classes
of planar distances.  At order fifteen this eliminates the exact cyclic
coloring with class sizes $(16,16,16,16,16,16,9)$ that survives all
one-point deletion tests.  This is a joint-realizability obstruction, not a
resolution of either unrestricted question in Erd\H{o}s Problem 132.
\end{abstract}

\section{The cyclic rigidity lemma}

Indices below are taken modulo $m$.  A regular star polygon means the orbit
of one point under a rotation through $2\pi k/m$, where $\gcd(k,m)=1$;
the label order need not be the convex circular order.

\begin{theorem}[Three-step rigidity]\label{thm:rigidity}
Let $m\ge5$ be odd, and let $p_0,\ldots,p_{m-1}$ be distinct points of
$\R^2$.  Suppose that, for each $h\in\{1,2,3\}$, the quantity
\[
  |p_{i+h}-p_i|
\]
is independent of $i$.  Then the points are the vertices of a regular star
$m$-gon in label order.  More precisely, after identifying $\R^2$ with
$\mathbb C$, there are $a\ne0$, $b\in\mathbb C$, and a primitive $m$th
root of unity $z$ such that
\[
  p_i=b+a\frac{1-z^i}{1-z}\qquad(0\le i<m).
\]
\end{theorem}

\begin{proof}
Put $e_i=p_{i+1}-p_i$.  The step-one hypothesis gives $|e_i|=s>0$ for
all $i$.  The step-two hypothesis gives
\[
 |e_i+e_{i+1}|^2=2s^2+2e_i\mathbin{\cdot}e_{i+1},
\]
so $e_i\cdot e_{i+1}=c$ is independent of $i$.

The two equality cases of Cauchy--Schwarz are impossible.  If $c=s^2$,
then all $e_i$ are equal, contradicting $\sum_i e_i=0$.  If $c=-s^2$,
then $e_{i+1}=-e_i$ around an odd cycle, again a contradiction.  Hence
there is a unique $\theta\in(0,\pi)$ such that
\[
 c=s^2\cos\theta.
\]
For each $i$, let $\sigma_i\in\{1,-1\}$ record whether the oriented angle
from $e_i$ to $e_{i+1}$ is $\theta$ or $-\theta$.

The step-three hypothesis and
\[
 |e_i+e_{i+1}+e_{i+2}|^2
 =3s^2+4c+2e_i\mathbin{\cdot}e_{i+2}
\]
show that $e_i\cdot e_{i+2}$ is independent of $i$.  On the other hand,
\[
 e_i\mathbin{\cdot}e_{i+2}
 =s^2\cos\bigl((\sigma_i+\sigma_{i+1})\theta\bigr).
\]
This is $s^2\cos(2\theta)$ when consecutive signs agree and $s^2$ when
they disagree.  Since $0<\theta<\pi$, these two values are different.
Thus either every consecutive sign pair agrees or every consecutive sign
pair disagrees.  The second alternative would alternate the signs around an
odd cycle and is impossible.  Therefore all $\sigma_i$ have one common
value.

Identify the plane with $\mathbb C$ and set $z=e^{\mathrm i\sigma_0\theta}$.
Then $e_i=e_0z^i$.  Closure gives
\[
 0=\sum_{i=0}^{m-1}e_i=e_0\frac{1-z^m}{1-z}.
\]
Here $z\ne1$, so $z^m=1$.  If the order of $z$ were a proper divisor
$d<m$, then
\[
 p_d-p_0=e_0\frac{1-z^d}{1-z}=0,
\]
contrary to distinctness.  Hence $z$ is primitive, and summing the edge
vectors gives the displayed formula.
\end{proof}

\begin{corollary}[Distinct cyclic chord classes]\label{cor:chords}
Under the hypotheses of Theorem~\ref{thm:rigidity}, the lengths belonging
to circular label differences $1,\ldots,(m-1)/2$ are pairwise distinct.
\end{corollary}

\begin{proof}
Write $z=e^{2\pi\mathrm i k/m}$ with $\gcd(k,m)=1$.  The squared length
at difference $h$ is a common positive factor times
\[
 |1-z^h|^2=2-2\cos(2\pi kh/m).
\]
Equality for differences $h$ and $j$ is equivalent to
$kh\equiv\pm kj\pmod m$, hence to $h\equiv\pm j\pmod m$.  In the stated
half-range this forces $h=j$.
\end{proof}

\section{The abstract one-rare-color family}

Let $n=2t+1$ be odd.  Label the vertices of $K_n$ by $\Z/n\Z$, and give
the edge $\{a,b\}$ its circular-difference color
\[
  j=\min\{|a-b|,n-|a-b|\}\in\{1,\ldots,t\}.
\]
Every color initially has $n$ edges.  Move $t-1$ distinct edges of color
$t$, one into each color $1,\ldots,t-1$.  The resulting color-class sizes
are
\[
  \underbrace{n+1,\ldots,n+1}_{t-1\text{ classes}},\quad t+2.
  \tag{1}
\]
They sum to $\binom n2$.  Thus, as abstract multiplicity data, (1) has
exactly one class of size at most $n$.

This family also defeats the immediate deletion induction.  Each lower
color retains its original 2-regular circular-difference graph and gains one
edge.  Deleting a vertex leaves either $n-2$ or $n-1$ edges of that color,
at most the new threshold $n-1$.  The top color is a subgraph of a
2-regular graph.  Hence every one-point deletion has many rare colors even
though the full abstract coloring has only one.

\begin{theorem}[Nonrealizability of the cyclic obstruction]
For every odd $n=2t+1\ge9$, the coloring above is not the partition of
$\binom n2$ pairs into Euclidean distance-equality classes of $n$ distinct
planar points.
\end{theorem}

\begin{proof}
Assume a realization $p_0,\ldots,p_{n-1}$ exists.  None of the original
edges in circular-difference colors $1,2,3$ was removed, so the step-one,
step-two, and step-three lengths are respectively constant.  (Extra moved
edges in those colors impose additional equalities but remove none of these
base equalities.)  Theorem~\ref{thm:rigidity} and
Corollary~\ref{cor:chords} show that circular differences
$1,\ldots,t$ have pairwise distinct lengths.  But every moved edge has
circular difference $t$ and was placed in one of the lower colors, requiring
its length to equal a lower-difference length.  This is a contradiction.
\end{proof}

\section{Order fifteen and scope}

At $n=15$ we have $t=7$, and (1) is precisely
\[
 (16,16,16,16,16,16,9).
\]
The last color can be chosen as a 15-cycle with six edges removed, hence a
linear forest; individually, such a graph can occur as a planar diameter
graph.  The theorem proves that the full seven-color system nevertheless
has no simultaneous Euclidean realization.  The obstruction therefore uses
joint metric compatibility rather than pair counting, deletion criticality,
or a property of the proposed diameter graph in isolation.

For Erd\H{o}s Problem 132 this removes one natural extremal family, including
the exact cyclic candidate isolated in the preceding continuation.  It does
not classify all seven-distance colorings on fifteen vertices, prove
$R\ge2$ at order fifteen, or imply that the unrestricted minimum of $R$
tends to infinity.  The remaining finite problem requires a reduction from
arbitrary candidate colorings to rigid joint patterns, while the asymptotic
problem still requires control of distances between differently oriented
triangular components and the large outer convex layers.

\section*{Reproducibility note}

The companion script \texttt{verify-cyclic-rigidity.py} checks the finite
edge counts, all one-vertex deletion counts, the order-fifteen profile, and
the modular distinct-chord implication for primitive roots.  The analytic
turning-angle argument in Theorem~\ref{thm:rigidity} is the load-bearing
general proof; the script is an exact consistency check, not a substitute for
that proof.

\begin{thebibliography}{9}
\bibitem{problem132}
T.~F. Bloom, \emph{Erd\H{o}s Problem 132},
\url{https://www.erdosproblems.com/132}.

\bibitem{additive}
B. Stone, \emph{Arithmetic-progression covers and rare planar distances},
repository manuscript,
\url{https://github.com/blakestone-x/math/tree/main/Erdos-132-convex-lattice}.
\end{thebibliography}

\end{document}
